\documentclass[11pt,a4paper]{article}

% -------------------- PACKAGES --------------------
\usepackage[margin=1in]{geometry}
\usepackage{amsmath,amssymb}
\usepackage{graphicx}
\usepackage{booktabs}
\usepackage{xcolor}
\usepackage{hyperref}
\usepackage{setspace}
\usepackage{fancyhdr}
\usepackage{lastpage}
\usepackage{authblk}
\usepackage{tcolorbox}
\usepackage{tocloft}

% -------------------- PAGE STYLE --------------------
\pagestyle{fancy}
\fancyhf{}
\cfoot{\thepage\ / \pageref{LastPage}}
\renewcommand{\headrulewidth}{0pt}
 

% -------------------- COMMENT BOX STYLE --------------------
\newtcolorbox{reviewbox}{
  colback=green!8!white,
  colframe=green!60!black,
  boxrule=0.7pt,
  arc=2pt,
  left=8pt,
  right=8pt,
  top=6pt,
  bottom=6pt
}

\newtcolorbox{updatebox}{
  colback=blue!6!white,
  colframe=blue!60!black,
  boxrule=0.7pt,
  arc=2pt,
  left=8pt,
  right=8pt,
  top=6pt,
  bottom=6pt
}

% -------------------- HIGHLIGHT --------------------
\newcommand{\rev}[1]{\colorbox{yellow}{#1}}

% -------------------- TITLE --------------------
\title{2nd Round Responses to Reviewers’ Comments for Manuscript\\[0.8em]
\textbf{Implementation of Communication and Group Formation in the Agent Based Modelling of Forced Migrants: Mali Case Study}\\[0.8em]
Addressed Comments for Publication to\\
\textbf{The Journal of Artificial Societies and Social Simulation (JASSS)}
}

\author[1]{Edward King}
\author[1]{He-In Cheong}
\author[1]{Arnab Majumdar}
\affil[1]{Department of Civil and Environmental Engineering, Imperial College London\\
South Kensington Campus, London SW7 2AZ}

\date{}

\begin{document}

% ======================================================
% ===================== PAGE 1 =========================
% ===================== TITLE PAGE =====================
% ======================================================

\maketitle
\vfill

\begin{center}
\textbf{Response Letter}
\end{center}

\newpage

% ======================================================
% ===================== PAGE 2 =========================
% ===================== CONTENTS =======================
% ======================================================

\setcounter{page}{1}
\renewcommand{\contentsname}{Contents}
\tableofcontents

\newpage

% ======================================================
% ===================== PAGE 3 =========================
% ===================== PREFACE ========================
% ======================================================

\section*{Preface}
\addcontentsline{toc}{section}{Preface}

[INSERT FINAL BRIEF HERE]

\vspace{1.2em}

\textbf{Note:} To enhance the legibility of this response letter, all the editor’s and reviewers’ comments are typeset in boxed format. The text inside the light green boxes represents the corresponding original reviewer comments. The updated sections are provided in a blue box. 

\newpage

% ====================================================
\section*{Response to Reviewer 1}
% ====================================================
\addcontentsline{toc}{section}{Response to Reviewer 1}

\subsection*{Abstract}

\begin{reviewbox}
Remove the explicit mention of the algorithm for network distances, it is not innovative enough that you would like to mention it at the beginning of the paper. Just say this, without the brackets:
\end{reviewbox}

\textbf{Response:}

Agreed. This has now been removed.

\vspace{0.5cm}

\begin{updatebox}
The model 
incorporates both physical and network-based path distances between nodes, and uniquely 
integrates dynamic border status changes over time, capturing real-world movement constraints and resistance across 
international borders. 
\end{updatebox}

\vspace{0.5cm}

\subsection*{Introduction}

\subsubsection*{Section 1.3}

\begin{reviewbox}
Drop the word "unethical", I think what you wanted to say is that "engineering population changes by using forced migration would constitute unethical behaviour/policies"; however, that section does not claim that either "you, the researchers" or "governments" do it, but simply that forced migration affects demography. This argument is uncontroversial, and does not have an ethical dimension.

The model incorporates both physical and network-based path distances between
nodes, and uniquely integrates dynamic border status changes over time, capturing
real-world movement constraints and resistance across international borders.
\end{reviewbox}

\textbf{Response:}

That was correct in what we were trying to say. However, as you pointed out it is redundant in this paragraph (and research). It has thus been removed.

\vspace{0.5cm}


\begin{updatebox}

The dynamics of forced migration are uncertain and highly complex, they are considered the most un-predictable
 drivers of population change    \citep{NRC2000}.

\end{updatebox}

\vspace{0.5cm}

\subsubsection*{1.9}

\begin{reviewbox}
Maybe replace "communication networks" with "social networks" because it is ambiguous, it may indicate "telecommunications" the way it is expressed.
Also "expalin" to "explain".
\end{reviewbox}

\textbf{Response:}

We agree with the ambiguity in this statement. It as been updated as you have suggested. The spelling mistake has also been corrected.

\vspace{0.5cm}

\begin{updatebox}

    The aim of this research is to enhance the accuracy of an ABM for forecasting the movement of forced migrants by 
integrating group dynamics and communication networks.

    \medbreak

    \textbf{Also in Appendix A:} The range of interaction is therefore defined both by explicit social networks and by emergent group 
membership. These assumptions reflect empirical observations that displaced populations rely heavily on word-of-mouth 
and kinship ties to navigate uncertain environments.

    \medbreak

    For brevity, the Supplementary Notes will be used to explain concepts and 
assumptions referenced in the main text in higher detail.

\end{updatebox}

\vspace{0.5cm}

\subsection*{Literature Review}

\subsubsection*{Section 2.1}

\begin{reviewbox}
Replace this:

"Its validation focuses on final outcomes, overlooking key decision-making processes, which limits its ability to capture
the full complexity of forced migrant movements (Suleimenova et al. 2017)."

With this:

"The model, insofar as it does not consider the interaction between agents and groups, is missing a key component of the real-world behaviour of agents, whose choices are affected by the choices of other agents".
\end{reviewbox}

\textbf{Response:}

We thank the reviewer for the clarity and phrasing, improving the previous phrase which can appear
ambiguous.

\vspace{0.5cm}

\begin{updatebox}
The model, 
insofar as it does not consider the interaction between agents and groups, is
missing a key component of the real-world behaviour of agents, whose choices are affected
by the choices of other agents.
\end{updatebox}

\vspace{0.5cm}

\subsection*{Methodology}

\subsubsection*{Section 3.3}

\begin{reviewbox}
You wrote a footnote that defines what is an IDP and what is a refugee. Why don't you use it in section 3.3.? When you talk about "refugees hosted in Mali", do you mean Malian citizens or foreigners? If the former, replace "refugees" with "IDPs".
\end{reviewbox}

\textbf{Response:}

Thank you for highlighting this. I've now applied the distinction consistently in Section 3.3 and clarified that “refugees hosted in Mali” refers to foreign nationals, in line with the definitions provided in the footnote.

\vspace{0.5cm}

\begin{updatebox}

At the heart of the Central Sahel, Mali faces multifaceted challenges including inadequate security, conflict, and 
climatic shocks. According to \citet{unhcr_refugee_statistics_2025} statistics, in 2012 (the year of the case study) 
Mali hosted approximately 14,000 foreign refugees while about 150,000 Malian nationals were displaced abroad as refugees 
\citep{unhcr_mali}. By 2024, these figures had risen to around 136,000 foreign refugees hosted in Mali and 354,000 Malian 
refugees residing in other countries \citep{unhcr_mali}. Beyond displacement, Mali grapples with severe environmental 
hazards such as rising temperatures, projected to increase by 2.5°C by 2080 \citep{76-wachiaya2021warming}, as well as 
risks of deforestation, desertification, and soil erosion that compromise water availability. Human trafficking also 
remains prevalent along migration routes, with digital technologies playing a dual role as both a risk factor and a 
resource for safety information \citep{UNHCR2022Trafficking}. Regionally, ongoing conflicts in the Central Sahel have 
led to more than 410,000 refugees and approximately 2.1 million internally displaced persons 
\citep{80-cheshirkov2022decade}.

\end{updatebox}

\vspace{0.5cm}

\subsubsection*{Section 3.4}

\begin{reviewbox}
When you talk about the variable "wealth distribution" it is not clear what you mean and how you built it. What I understand is that you took the measurement for Mali in the World Inequality Dataset, then you found from Oxfam the cumulative capital held by the 0.1\% of the top distribution of income, and then you fit this falue into the distribution associated with the measure of inequality in the World Inequality Dataset? Please expand and clarify how exactly this variable is built, it remains unclear
\end{reviewbox}

\textbf{Response:}
Thank you for this comment. I have revised the manuscript to clarify the distinction between data sources and model implementation. Section 3.4 now only describes the datasets used (including the World Inequality Dataset), while the detailed construction of the wealth variable has been moved to Section 4.5, where the full functional form and implementation within the ABM are explicitly defined. This separation improves clarity and reproducibility.

\vspace{0.5cm}

\begin{updatebox}
GRASP is an ABM designed to predict immediate forced migrant flows after sudden conflicts, utilising empirical evidence on the grouping and communication of forced migrants. Although applied to Mali 2012 with UNHCR micro-data \citep{He-In}, GRASP can be extended to other cases using a combination of qualitative and quantitative data, including conflict event data \citep{243-acled2024data}, demographic structure (age and gender distributions), settlement and camp population data, and income inequality information derived from the World Inequality Dataset.

\medbreak

\textbf{Section 4.5: } In the model, wealth influences agents’ transport speed and destination selection. Wealth is assigned using a percentile-based transformation of the form:
\begin{equation}
W(p) = 100000 \cdot (100 - p)^{-0.25} - 30000,
\end{equation}
where \(p \in [0,100]\) denotes the agent’s wealth percentile. Age affects both mobility and dependent classification (relevant to group formation), while gender is not currently used but remains a possible future extension.
\end{updatebox}

\vspace{0.5cm}


\begin{reviewbox}

Sudan is ok, Moldova and Myanmar are not good comparisons. Add a sentence that states something like "For the purpose of this particular paper, which focuses more on the methodology than on the exact replication of migrant distributions for Mali, we assume that migrants in Mali behave similarly to those from Moldova or Myanmar".

\end{reviewbox}

\textbf{Response:}

Thank you for the suggestion. I've added a clarification acknowledging this limitation and explicitly stated the behavioural assumption underlying the use of Moldova and Myanmar data.

\vspace{0.5cm}

\begin{updatebox}

    Micro-data were extracted from UNHCR case studies on Sudan, Moldova, and Myanmar by \citet{He-In}. Household size 
distributions, though not directly related to Mali, are comparable to Sudan due to similar mean household sizes 
\citep{dhs_sr261,ESRI_AverageHousehold}. Moldova and Myanmar were the only available sources for data on forced migrant 
travel groups and their sources of information. For the purpose of this paper, which focuses on the methodological 
framework rather than the exact replication of migrant distributions in Mali, we assume that migrants in Mali behave 
similarly to those observed in the Moldova and Myanmar case studies.

\end{updatebox}

\vspace{0.5cm}

\subsubsection*{Section 3.11}

\begin{reviewbox}
The confusion matrix seems superfluous, and it occupies half of a page, I would personally remove it but it's your choice. Virtually all readers of the paper will understand the terminology contained in it regardless of your picture.
\end{reviewbox}

\textbf{Response:}

We agree that it is safe to assume the the typical reader will be familiar with the confusion matrix.
Further explanation has also already been provided in the supplementary note. Therefore, we have removed
the figure and updated the paragraph.

\vspace{0.5cm}

\begin{updatebox}
Estimating the number of forced migrants can be viewed as a form of binary classification (see Supplementary note 5). 
For this reason, in our analysis, we will focus on precision and accuracy. 
Precision is the ratio of correctly classified positive outcomes to total 
predicted positive outcomes, while accuracy is the total number of correctly predicted outcomes \citep{confusion}.
\end{updatebox}

\vspace{0.5cm}

\subsubsection*{Section 3.13 and Section 3.14}

\begin{reviewbox}
Sections 3.13 and 3.14 seem unnecessary. You can simply indicate the measure for accuracy that you chose.
\end{reviewbox}

\textbf{Response:}

??? ARNAB Agree we can remove 3.14 but is it not worth explaining (even if shorter) why precision is more
important in this case than accuracy.

\vspace{0.5cm}

\begin{updatebox}

\end{updatebox}

\vspace{0.5cm}

\subsection*{Agent Based Models}

\subsubsection*{Section 4.9}

\begin{reviewbox}
Agents who are just born are unlikely to be able to travel alone, but to my understanding you have not modeled "this agent follows exaclty the same trajectory as the agent called "parent"", so births in your model essentially indicate "an agent has come to age and can in principle move autonomously, if selected to do so; otherwise, they travel with their family".
\end{reviewbox}

\textbf{Response:}
Thank you for this comment. I have revised the text to explicitly clarify the treatment of newborn agents in the model. In particular, I now state that newborns are immediately assigned to a family group and do not make independent migration decisions, with their movement fully determined by group dynamics. I have also added clarification that the absence of an age-based transition to autonomy is a current limitation of the model.

\vspace{0.5cm}

\begin{updatebox}
Family sizes are defined according to micro-data from Sudan\footnote{Ref}. Families are formed at the start of the simulation, with the possibility of increasing in size through births and decreasing through deaths. Newborn agents are immediately assigned to a family group and do not make independent migration decisions; instead, their movement is fully determined by the behaviour of the group to which they belong. While births are included primarily to support long-term simulations and group cohesion, their impact is minimal over shorter time frames. A limitation of the current model is that there is no age-based transition to independent decision-making. Deaths, however, play a more active role—especially in conflict zones—where they influence both agent survival and the dissemination of risk information to other agents.
\end{updatebox}

\vspace{0.5cm}

\subsubsection*{Section 4.20}

\begin{reviewbox}
In section 4.20, you could consider adding a point at the end that states "in principle, if real-world trajectories sampled at daily frequency were available, these beta parameters could be learnt". For you personally, not for the paper, I also recommend thinking whether a form of reinforcement learning for the model's parameters can be established:

- initialise beta weight randomly or by picking them
- compute error/accuracy after simulating everything
- compute the effect that the weights have on the error and do backpropagation
- adjust the weights with some small steps in the direction of the steepest ascent/descent
- run the simulation again.

This is one approach by means of which you can attempt to crack the parametrisation of the model via gradient descent, and would in the meanwhile also generate sintetic trajectories that you could use for other purposes. This is however something to think about for another paper.
\end{reviewbox}

\textbf{Response:}

Thank you for this excellent suggestion. This is definitely a very promising avenue for further research. I have added a note on the potential to learn the $\beta$ parameters from trajectory data; however, implementing such an approach would substantially extend the scope and complexity of the current simulation, so it is left for future work.

\vspace{0.5cm}

\begin{updatebox}

In principle, if real-world trajectories sampled at a daily frequency were available, these $\beta$ parameters could be learned directly from data.
In addition, further clarity has been provided in the Future Research Section.

\textbf{Added to Future Research} Future developments in computational intelligence methods, such as data mining and 
reinforcement learning, could reflect 
realistic agent behaviour, provided they are justified and not driven by trends. \textbf{In particular, reinforcement learning or 
gradient-based approaches could be explored to calibrate behavioural parameters (e.g. $\beta$ weights) from observed or 
synthetic trajectory data.} The goal is to balance the improvement 
in granularity with the increase in computational complexity. 
\end{updatebox}

\vspace{0.5cm}

\subsection*{Results}

\subsubsection*{Figure 6}

\begin{reviewbox}

Here are the images obtained, which do not match those presented in the paper:

\href{https://pasteboard.co/LzLPAUiGF1m9.png}{Image 1}
\href{https://pasteboard.co/KUbKqpYE6Ia1.png}{Image 2}

Additionally, the file \texttt{Country\_split\_refugees\_w0\_roulette.csv} is not present in the repository. 
While this would typically prevent replication, this can still be resolved prior to resubmission.

To ensure reproducibility, the repository should allow an interested reader to clone it, run all scripts, and reproduce the figures exactly as presented in the paper. Otherwise, the reliability of the results is significantly reduced.
\end{reviewbox}

\textbf{Response:}


\vspace{0.5cm}

\subsubsection*{Section 5.5}

\begin{reviewbox}
Section 5.5. has a problem, fix it before final release, let us pretend nobody noticed it.
\end{reviewbox}

\textbf{Response:}

I have removed this editors note that should not have been submitted. There was also another
comment associated with Section 6.9 that has been addressed. In both cases, the appropriate
reference link has been added.

\vspace{0.5cm}

\subsubsection*{Section 5.13}

\begin{reviewbox}
section 5.13 is too long, please cut it in half. I would remove everything after footnote 12.
\end{reviewbox}

\textbf{Response:}

We agree that the paragraph is too lengthy and that beyond footnote 12, there is very little
 added analysis. This part has now been removed.

\vspace{0.5cm}

\begin{updatebox}
Fig. \ref{flee0} compares Flee and GRASP for Mentao and Abala (with all 7 comparable camps in Fig. S7 of the 
Supplementary Information). When scaled in the same fashion as Flee, it becomes apparent that GRASP is unable to predict 
forced migrant flow with a higher degree of precision than Flee. Disregarding the transition peaks, GRASP accurately 
identifies peaks by the date and produces a ballpark estimate of the number of refugees. Large gradients are also 
detected, which are crucial for humanitarian support as they indicate when immediate aid is required for a camp. 
Regarding the characteristic shape of the forced migrant flows, Mentao, Mbera, Abala, Mangaize, and Tabareybarey all 
indicate areas of population stability and areas of population increase/decrease to a higher degree of accuracy than 
Flee, even if they arrive at the wrong population value\footnote{This is particularly important for humanitarian 
operations if the refugee registrations are not reliable.}. 
\end{updatebox}

\vspace{0.5cm}

\subsubsection*{Figure 12}

\begin{reviewbox}
Ok I understand you don't want to change ylim in the top figure, but you really cannot keep half the chart empty in the bottom figure. Please fix it by rescaling it and limit to y=1.25.
\end{reviewbox}

\textbf{Response:}

\vspace{0.5cm}

\begin{updatebox}

\end{updatebox}

\vspace{0.5cm}

\subsection*{Discussion}

\subsubsection*{Section 6.2}

\begin{reviewbox}
" leading to unrealistic simulation" is repeated
\end{reviewbox}

\textbf{Response:}

We have corrected this. We thank the reviewer for spotting this mistake.

\vspace{0.5cm}

\begin{updatebox}
GRASP addresses limitations in current forced migration ABMs like the Flee 1.0 model, which fail to account for forced 
migrant experiences and micro-processes, leading to simulations with unverified 
fidelity. This research assesses the effectiveness of incorporating migrant topology (the interactive element of agents 
between themselves and their environment), focusing on grouping and communication, to improve forecasting in forced 
migration ABMs. These aspects were chosen based on advancements in literature 
\citep{84-hinsch2019rumours,163-bijak2021routes,121-collins2016agent-based}. The dynamic weighting of nodes, based on 
evolving conflict data and border accessibility, also allows the model to reflect real-time constraints in migrant 
decision-making.
\end{updatebox}

\vspace{0.5cm}

\subsubsection*{Section 6.9}

\begin{reviewbox}
Fix the missing link if one is supposed to go there, otherwise remove it if it is a placeholder.
\end{reviewbox}

\textbf{Response:}

See above. This has alsop been fixed.

\vspace{0.5cm}

\subsubsection*{Section 6.12}

\begin{reviewbox}
I would like to understand how exactly expensive is it, and most readers will share the same question. Is it one laptop full time for one week? Or did you have to rent a cluster of GPUs for a year? You should insert in this section a clarification on the actual time it took you to run everything, because it seems to me that this type of model should not occupy that many computational resources.
\end{reviewbox}

\textbf{Response:}

\vspace{0.5cm}

\begin{updatebox}

\end{updatebox}

\vspace{0.5cm}

\subsubsection*{Section 6.14}

\begin{reviewbox}
Add a clarification that the reduction of 100 people to 1 agent may introduce distortions, when we consider that intra-family dynamics between agents are also present.
\end{reviewbox}

\textbf{Response:}

Thank you for this helpful observation. I've added a clarification noting that aggregating 100 individuals into a single agent may introduce distortions in intra-family dynamics.
This was included before cutting the report down for brevity but we agree that this point is important.
\vspace{0.5cm}

\begin{updatebox}

Additionally, this level of aggregation may introduce distortions in intra-family dynamics, as interactions that would occur within families are instead represented between agents.

\end{updatebox}

\vspace{0.5cm}

\subsubsection*{Section 6.18}

\begin{reviewbox}
You do have some type of wave-like behaviour, in the sense that the arrival of one agent at destination increases the likelihood that the same destination is chosen by another agent. Please add a statement of the form "however, our model can partially represent dependencies in the selection of migration trajectories by future agents on the basis of the choices by past agents insofar as the arrival at destination by one migrant can influence the future arrival to that same destination of other migrants".
\end{reviewbox}

\textbf{Response:}

Agreed. I've added a clarification noting that the model can partially capture such wave-like dependencies through destination-based interactions between agents.

\vspace{0.5cm}

\begin{updatebox}

However, the model can partially represent dependencies in the selection of migration trajectories, insofar as the arrival of one agent at a destination can increase the likelihood that subsequent agents select the same destination.

\end{updatebox}

\vspace{0.5cm}

\subsubsection*{Section 6.22}

\begin{reviewbox}
Remove this: "To improve scalability and deployment speed, future itera-
tions of GRASP could explore unsupervised methods, such as minimum spanning trees or distance-constrained
heuristics, to generate an initial network structure. These could offer a practical starting point in the immediate
aftermath of a crisis, before manual refinement introduces higher spatial and contextual accuracy"
\end{reviewbox}

\textbf{Response:}

We have removed this section. We agree and believe that this offers very little further analysis in an already 
lengthy Discussion section. It also has very little academic grounding with no continuation with the theme of the
 paper.

\vspace{0.5cm}

\begin{updatebox}

\end{updatebox}

\vspace{0.5cm}

\subsubsection*{Section 6.23}

\begin{reviewbox}
It would make sense to use a sentence of the form "the agent does not have global memory of the trajectory, but only a local memory of its immediate past; a future iteration of this study may include a long-term memory for the agent that is incrementally populated as the agent migrates". It seems like the concept of "global/local" is relevant and should be added there, both programmers and migration scholars will understand what it means.
Also drop the sentence "While this would increase memory and computational load", in this paper we are talking about comparatively little data, and the level of discretization you use make the problem computationally tractable.
\end{reviewbox}

\textbf{Response:}

I've incorporated the distinction between local and global memory and clarified the scope for extending agent memory in future work, while removing the note on computational cost.

\vspace{0.5cm}

\begin{updatebox}

The current model includes a limited form of memory, where agents retain and share knowledge of previously encountered 
conflict zones through a system of “no-go” nodes. This mechanism allows agents to avoid revisiting dangerous areas and 
contributes to more informed route planning. The agent does not have global memory of its full trajectory, but instead 
operates with a local memory of its immediate past interactions and exposures. 
However, this memory is specific to conflict exposure and does not extend 
to a broader recollection of the agent’s full path or personal experiences during migration. 
A future iteration of this study may incorporate a longer-term memory that is incrementally populated as the agent migrates. 
Future enhancements could 
build on this foundation by introducing fuller path-dependency, where agents remember the sequence of nodes traversed or 
assign qualitative experiences (e.g., hardship, delay) to locations. This would allow agents to refine decision-making 
based on accumulated experiences, such as avoiding inefficient or distressing routes, even in the absence of direct 
conflict.

\end{updatebox}

\vspace{0.5cm}

\subsubsection*{Section 7.3}

\begin{reviewbox}
Remove the bullet list and write in a single paragraph.
\end{reviewbox}

\textbf{Response:}

The bulleted list has been collapsed and slightly reworded for ease of reading.

\vspace{0.5cm}

\begin{updatebox}

Key findings indicate that even at a 1\% population scale, GRASP outperforms the Flee model in several ways. It is able to 
estimate the number of refugees in camps directly from conflict data with improved precision (a capability that Flee lacks 
without relying on refugee registration data). It also identifies sudden gradients, peaks, and periods of stability in UNHCR 
camp populations, and classifies not only refugees but also IDPs and returnees with greater precision. Additionally, GRASP 
utilises empirical data for the micro-processes involved in refugee movement where possible, rather than relying on assumed 
proxies, and provides a utility-based decision model that can be further refined through interviews with Malian residents and 
refugees. Finally, it is capable of estimating the number of non-camping refugees, including those seeking refuge in 
neighbouring settlements and those fleeing via airports.

\end{updatebox}

\vspace{0.5cm}

\subsection{Bibliography}

\begin{reviewbox}
This reference is wrong, you inserted probably "Article in journal", while it is a "Report": Network, G. O. (2017). Global opportunity report 2017. Tech. rep., DNV GLAS
I found it here, please fix it since otherwise it seems like Network is somebody's surname: https://www.dnv.com/publications/global-opportunity-report-2017-84879/
\end{reviewbox}

\textbf{Response:}

Thank you for catching this. I've corrected the reference to properly reflect it as a technical report rather than a journal article.

??? CHECK COMPILATION UPDATE WORKED

\vspace{0.5cm}

% ====================================================
\section*{Response to Reviewer 2}
% ====================================================
\addcontentsline{toc}{section}{Response to Reviewer 2}

\subsection*{Introduction}

\subsubsection*{Section 1.3}

\begin{reviewbox}
I don't know that the parenthetical "and unethical" is necessary here. I had to read it three times because it made the paragraph less clear.
\end{reviewbox}

\textbf{Response:}

We agree (as does reviewer 1). The bracketed part has now been removed.

\vspace{0.5cm}

\begin{updatebox}

The dynamics of forced migration are uncertain and highly complex, they are considered the most un-predictable
 drivers of population change    \citep{NRC2000}.
\end{updatebox}

\vspace{0.5cm}

\begin{reviewbox}
I'm unclear on the evidence to claim that "migration modelling itself remains underdeveloped" and that "models often struggle to produce reliable theoretical predictions." - these warrant citations.
\end{reviewbox}

\textbf{Response:}

Thank you for this comment. I have revised this statement to better reflect the existing literature and added supporting citations (Massey et al. 1993; Davenport et al. 2003; Van Hear et al. 2018; Cheong 2023) to substantiate the discussion of limitations in current migration modelling approaches.
These citations were used in a preliminary version of the report but have been reintroduced to add academic validity.

\vspace{0.5cm}

\begin{updatebox}
However, several studies highlight limitations in current migration modelling approaches, including weak theoretical foundations, methodological biases, and insufficient treatment of complex driver interactions (Massey et al. 1993; Davenport et al. 2003; Van Hear et al. 2018; Cheong 2023).
\end{updatebox}

\vspace{0.5cm}

\subsubsection*{Section 1.4}

\begin{reviewbox}
Citation for "new levels of granularity" is from 1990. Demographics modeling has progressed in granularity since that time, I believe.
\end{reviewbox}

\textbf{Response:}
Thank you, you make a good point regarding the age of the original citation. To address this, I have revised the text to reflect that increasing modelling granularity is an ongoing development and added a more recent reference (Bijak, 2022), which provides a contemporary overview of advances in migration modelling.

\vspace{0.5cm}

\begin{updatebox}
With the advancement in conceptualising population growth and decline mechanisms, improvement of modelling techniques, and rapid development of computational power, demographic forecasting has evolved to increasingly fine levels of granularity \citep{106-willekens1990demographic, bijak2022}.
\end{updatebox}

\vspace{0.5cm}

\subsubsection*{Section 1.5}

\begin{reviewbox}
I'm not sure I follow what this paragraph is trying to say about "researchers in host countries."
\end{reviewbox}

\textbf{Response:}
Thank you for this comment. I have revised the paragraph to clarify the role of researchers in host countries and to better link this discussion to the limitations it introduces in migration modelling.

\vspace{0.5cm}

\begin{updatebox}
These limitations can constrain the development of more comprehensive and representative migration models.
\end{updatebox}

\vspace{0.5cm}

\subsubsection*{Section 1.6}

\begin{reviewbox}
Literature on push/pull factors and claim that "theoretical understanding of the drivers behind migration remains weak" when citation dates to 1993.
\end{reviewbox}

\textbf{Response:}
Thank you, you raise a valid point regarding the age of the citation. I have revised the text to clarify that this limitation has been recognised for some time and remains an ongoing challenge, and I have added more recent references to reflect the current state of the literature.

\vspace{0.5cm}

\begin{updatebox}
The theoretical understanding of the drivers behind migration has long been recognised as limited, despite their significant impact across many countries \citep{227-massey1993theories}. More recent literature continues to highlight challenges in capturing the complexity and interaction of migration drivers. Historical models of migration have also been criticised for sample selection bias and insufficient attention to 'pull' factors (drivers in potential destination countries that attract forced migrants) \citep{145-davenport2003leave:}.
\end{updatebox}

\vspace{0.5cm}

\subsubsection*{Section 1.7}

\begin{reviewbox}
First sentence: "Existing models often overlook the role of group dynamics and communication among migration populations" -- no citations to back this up. Second sentence: "These elements are essential..." -- says who? for what? Later in the same paragraph "...current models struggle" -- what models? What specifically do they cover/not cover to identify that "gap in the literature"?
\end{reviewbox}

\textbf{Response:}
Thank you for this detailed comment. I have revised the paragraph to clarify what is meant by “existing models” and to better justify the importance of group dynamics and communication. Specifically, I now distinguish between individual-level and interaction-based modelling approaches and clarify how the omission of these elements limits model realism. Relevant clarifications have been added in the text.

\vspace{0.5cm}

\begin{updatebox}
Existing models often overlook the role of group dynamics and communication among migrant populations, with many migration models focusing primarily on individual-level decision-making or aggregate flows rather than interactions between agents. These elements are essential for understanding certain migration patterns, particularly when people travel together and share information to guide decisions about potential destinations, as highlighted in studies of migration networks and information flows. Without accounting for such interactions, these types of models struggle to capture the realities of these forms of forced migration, leaving a gap in the literature. At the same time, it is important to recognise that group-based movement is not universal. It is most commonly linked to acute-onset events, such as sudden outbreaks of violence or climate-related disasters. In other contexts, especially in some waves of displacement, solo migration is more prevalent. For instance, during the early stages of the 2015 refugee crisis in Lesvos, Greece, the first arrivals were often single men traveling alone, while family groups became more prominent in later waves as demographics shifted \citep{unhcr_sea_arrivals_2015,crawley2016unpacking}.
\end{updatebox}

\vspace{0.5cm}

\subsubsection*{Section 1.9}

\begin{reviewbox}
"...leads to improved predictive outcomes" of what? over what? in what specific dimensions of forced migration movement?
\end{reviewbox}

\textbf{Response:}
Thank you for this comment. I have revised the sentence to clarify what is meant by “improved predictive outcomes,” specifying the dimensions of migration being evaluated (e.g., camp population distributions, destination selection, and temporal dynamics) and the comparison to a baseline model (Flee), as well as explicitly incorporating both grouping and communication mechanisms.

\vspace{0.5cm}

\begin{updatebox}
Using the case study of Mali in 2012, the research seeks to demonstrate that a more granular understanding of the forced migrant experience leads to improved predictive outcomes, particularly in modelling camp population distributions, destination selection, and the temporal dynamics of displacement, through the incorporation of grouping mechanisms and communication processes, compared to baseline models such as Flee.
\end{updatebox}

\vspace{0.5cm}

\subsubsection*{Section 1.10}

\begin{reviewbox}
Typo in the word "explain"
\end{reviewbox}

\textbf{Response:}

This has now been corrected. We thank the reviewer for spotting the mistake.

\vspace{0.5cm}

\subsection*{Literature Review}

\subsubsection*{General Comment}

\begin{reviewbox}
The literature review still feels very weak to set up the importance of the model.
\end{reviewbox}

\textbf{Response:}

Thank you for this suggestion. We have strengthened the literature review by adding additional context on the importance of forecasting forced migration, the rationale for using agent-based models, and the limitations of existing approaches before introducing Flee and Metafora. We also expanded the discussion of the research gap to more clearly motivate the development of GRASP.

\vspace{0.5cm}

\begin{updatebox}

Input added parts

\end{updatebox}

\vspace{0.5cm}

\subsubsection*{Section 2.4}

\begin{reviewbox}
2.4 is the most compelling literature review paragraph in the first section. It best establishes the role of this particular model and is situated in more recent literature.
\end{reviewbox}

\textbf{Response:}

Thank you for this observation. We have used Section 2.4 as a model to strengthen the surrounding literature review sections.

\vspace{0.5cm}

\subsection*{Results}

\subsubsection*{Section 5.1}

\begin{reviewbox}
"While some deviations will exist between GRASP predictions and empirical camp populations, this does not undermine the value of the model" -- rather than stating this, can you provide a statistical summary that shows that variations are 1. on par or better than FLEE or other models? and/or 2. statistically insignificant? Normally we wouldn't expect models to be perfectly aligned with real-world data, but one would expect some kind of benchmark or statistical analysis to give a reference point for "good enough."
Similar comment in paragraph 5.3. The explanation feels defense relative to FLEE, rather than statistically convincing relative to GRASP's own V\&V benchmarks.
\end{reviewbox}

\textbf{Response:}

\vspace{0.5cm}

\begin{updatebox}

\end{updatebox}

\vspace{0.5cm}

\subsubsection*{Section 5.6}

\begin{reviewbox}
About the limitations re: utility function reliance -- what impact does this have? The explanation about lengthy run times falls a bit flat rather than focusing on what we can know from this exercise.
\end{reviewbox}

\textbf{Response:}
Thank you for this comment. I have revised this section to focus on the implications of the utility function as a modelling limitation, rather than computational constraints. In particular, I now clarify how parameter sensitivity in the utility function affects destination and decision-level predictions, and why this motivates further qualitative calibration of micro-level behavioural assumptions.
\vspace{0.5cm}

\begin{updatebox}
The main limitation of the GRASP model lies in its reliance on a calibrated utility function for decision-making. While parameter optimisation was performed at higher \textit{frac} values due to computational constraints, this introduces uncertainty when the model is scaled to full-resolution simulations, as parameter sensitivity may vary with population size. As a result, predictions related to destination choice and migration decision pathways are particularly sensitive to utility specification, which may affect fine-grained accuracy in camp selection outcomes (Fig. \ref{emp0}). It is therefore important not to rely solely on aggregate model output for validation, as \citet{He-In} note that this approach is insufficient. Instead, evaluating the underlying micro-processes is essential, motivating the incorporation of qualitative data from forced migrants to better inform and constrain the utility-based decision mechanism.
\end{updatebox}

\vspace{0.5cm}


\end{document}

